% Generated draft: requires mathematical and source review.
\documentclass[11pt]{article}
\usepackage[margin=25mm]{geometry}
\usepackage{fontspec}
\setmainfont{Latin Modern Roman}
\usepackage{amsmath,amssymb,mathtools,unicode-math}
\setmathfont{Latin Modern Math}
\usepackage{enumitem,needspace,xurl,hyperref}
\setlength{\parindent}{0pt}
\setlength{\parskip}{5pt}
\setlength{\emergencystretch}{3em}
\newcommand{\N}{\mathbb{N}}
\newcommand{\Z}{\mathbb{Z}}
\newcommand{\Q}{\mathbb{Q}}
\newcommand{\R}{\mathbb{R}}
\newcommand{\C}{\mathbb{C}}
\newcommand{\F}{\mathbb{F}}
\newcommand{\A}{\mathbb{A}}
\newcommand{\PP}{\mathbb{P}}
\newcommand{\E}{\mathbb{E}}
\newcommand{\eps}{\varepsilon}
\newcommand{\dd}{\,\mathrm{d}}
\begin{document}
\section*{UnsolvedMath ID: 30005901}
\textbf{Source problem code:} catalogue entry
\subsection{Definitions and mathematical statement}
\textbf{Scope and unresolved formulation.} The supplied record concerns epigenetics in complex living systems. It does not define a mathematical state space, time domain, evolution law, class of observations, admissible bioinformatic methods, or criterion for understanding such systems. Consequently, it supplies no specified mathematical objects over which hypotheses or conclusions can be quantified, and no precise theorem or conjecture can be extracted without additional assumptions.

\textbf{Source question.} The original question is: ``What is the next step in the field of bioinformatics to explore and understand epigenetics in complex life systems?'' The original statement appends the source locator \emph{Mini-Workshop: Mathematics of Entropic AI in the Natural Sciences}, page 1065.

The registry formulation is: ``Which bioinformatic and mathematical problems are central to understanding epigenetics as a complex dynamical system?'' This is an open research-direction question. The original wording asks for a next step in bioinformatics; the registry wording explicitly introduces mathematical problems and a dynamical-system interpretation. Their precise intended relationship is unresolved by the supplied evidence.

\textbf{Statement-only target.} Identify and formulate central bioinformatic and mathematical research problems for exploring and understanding epigenetics in complex living systems, with the dynamical-system interpretation appearing in the registry formulation. The source specifies no unique output, ranking criterion, hypotheses, boundary conditions, or success criterion. Any particular dynamical model, inference problem, or causal identification target would therefore require additional choices not supplied here.

The supplied dated curation mentions single-cell and multimodal epigenomic integration, spatial context, and causal interpretation as continuing challenges. These are contextual research directions, not defined subproblems or mandatory components of a mathematical statement. The supplied label \emph{open} is a curation status, not an independently verified assertion about a precisely formulated conjecture.
\subsection{Short English statement}
The record asks which research directions could advance understanding of epigenetics in complex living systems. Its registry version emphasizes mathematical problems and dynamics, but supplies no model or testable mathematical target. A precise conjecture requires further specification; the draft preserves the question and its unresolved scope.
\subsection{Sources}
\begin{itemize}
\item Local UnsolvedMath registry: supplied complete record, permanent ID \(30005901\), exact source problem code \texttt{catalogue entry}, title \emph{Bioinformatic Frontiers in Complex-System Epigenetics}; collection: \emph{Oberwolfach Reports Open Problems}. The record gives proposed year 2024 and names no individual proposer.
\item Supplied report citation: \emph{Mini-Workshop: Mathematics of Entropic AI in the Natural Sciences} (2024), Oberwolfach Reports, pp.\ 1041--1068, DOI 10.4171/owr/2024/18. The original statement identifies page 1065. \url{https://doi.org/10.4171/owr/2024/18}. The supplied background describes Oberwolfach Reports as documenting mathematical workshops and their open-problem discussions; the report has not been independently consulted for this draft.
\item Supplied literature source: ``Argelaguet et al., single-cell multi-omics.'' \url{https://www.nature.com/articles/s41580-023-00615-w}.
\item Supplied literature source: ``Single-cell and multimodal omics review.'' \url{https://www.nature.com/articles/s12276-024-01186-2}.
\item Literature assessment dated exactly \texttt{2026-08-22}: the supplied curation labels the question open, describes it as a broad future-direction question with no unique established solution, and identifies single-cell and multimodal epigenomic integration, spatial context, and causal interpretation as continuing challenges. This is supplied dated curation derived from the OWR record, not an independently verified literature assessment. The supplied background requests independent verification of time-sensitive status and cited claims.
\item The registry labels the statement \texttt{corrected\_verified} and states that only formula delimiters, notation, grammar, and typography were repaired while mathematical content and scope were preserved. This supplied verification label does not establish a precise mathematical formulation or resolve the wording difference recorded above.
\end{itemize}
\end{document}
